% 六册共用的导言区。各册 main.tex 在 \documentclass 之后 \input 本文件，
% 再追加本册专用的宏包与命令。改动这里会影响全部 6 册，改完请逐册编译检查。
%
% 宏包加载顺序会影响结果（如 hyperref 应较晚加载），调整顺序前请先确认。

\usepackage[margin=2.5cm]{geometry}

% 脚注用带圈数字 ①②③……，每页重新编号
\usepackage{pifont}
\usepackage[perpage,symbol*]{footmisc}
\DefineFNsymbols{circled}{{\ding{192}}{\ding{193}}{\ding{194}}
{\ding{195}}{\ding{196}}{\ding{197}}{\ding{198}}{\ding{199}}{\ding{200}}{\ding{201}}}
\setfnsymbol{circled}

\usepackage{ulem}

\usepackage{amsmath,amsfonts,mathrsfs,amssymb}
\usepackage{graphicx}

\usepackage[font=bf,labelfont=bf,labelsep=quad]{caption}

\usepackage{tikz}

\usepackage{ntheorem}
\theoremseparator{\;}

\usepackage{blkarray}
\usepackage{bm}
\usepackage[colorlinks=true, linkcolor=black]{hyperref}

\usepackage{enumerate}

% 选做题：用 \staritem 代替 \item，编号照常递增，题号左侧加“*”（同原书）
\makeatletter
\newcommand{\staritem}{%
  \refstepcounter{\@enumctr}%
  \item[\llap{*}\csname label\@enumctr\endcsname]}
\makeatother

\theoremstyle{plain}
\theoremheaderfont{\normalfont\bfseries}
\theorembodyfont{\normalfont}

\usepackage[framemethod=tikz]{mdframed}

% ---- 例题、解答等环境 ----
\newtheorem{example}{\bf 例}[chapter]
\newenvironment{solution}{\noindent {\bf 解：}}{}
\newenvironment{analyze}{\noindent {\bf 分析：}}{}
\newenvironment{rmk}{\noindent {\bf 注意：}}{}
\newenvironment{note}{\noindent {\bf 说明：}}{}

\renewcommand{\proofname}{\bf 证明：}
\newenvironment{proof}{{\noindent \bf 证明：}}{}%{\hfill $\square$\par}

% ---- 常用数学命令 ----
\newcommand{\E}{\mathbb{E}}
\renewcommand{\Pr}{\mathbb{P}}
\newcommand{\EP}{\mathbb{E}^{\mathbb{P}}}
\newcommand{\EQ}{\mathbb{E}^{\mathbb{Q}}}
\newcommand{\dif}{\,{\rm d}}
\newcommand{\Var}{{\rm Var}}
\newcommand{\Cov}{{\rm Cov}}
\newcommand{\x}{\times}

% ---- “练习”框与带标题的框 ----
\usepackage{tcolorbox}
\tcbuselibrary{breakable}
\tcbuselibrary{most}

\newtcolorbox{ex}[1][]
  {colback = white, colframe = cyan!75!black, fonttitle = \bfseries,
    colbacktitle = cyan!85!black, enhanced,
    attach boxed title to top center={yshift=-2mm},breakable,
    title=练习, #1}

\newtcolorbox{blk}[2][]
  {colback = white, colframe = magenta!75!black, fonttitle = \bfseries,
    colbacktitle = magenta!85!black, enhanced,
    attach boxed title to top left={xshift=5mm, yshift=-2mm},breakable,
    title=#2, #1}

% ---- 章节编号：第一节 / 一、/（一） ----
\setcounter{tocdepth}{2}
\setcounter{secnumdepth}{3}

\ctexset {
section = {
	name = {第,节},
 	number = \chinese{section}},
subsection = {
	name = {,、\hspace{-1em}},
	number = \chinese{subsection}
},
subsubsection = {
	name = {（,）\hspace{-1em}},
	number = \chinese{subsubsection},
}
}

\renewcommand{\contentsname}{目~~录}

\newcommand{\poly}{\polynomial[reciprocal]}

\usepackage{mathtools}

\setlength{\abovecaptionskip}{0.cm}
\setlength{\belowcaptionskip}{-0.cm}

\usetikzlibrary{decorations.pathmorphing, patterns}
\usetikzlibrary{calc, patterns, decorations.markings}
\usetikzlibrary{positioning, snakes}

\usepackage{yhmath}   % \wideparen
\usepackage{longdivision}
\usepackage{polynom}
\usepackage{polynomial}
\usepackage{cancel}

\renewcommand{\frac}{\dfrac}
\newcommand{\oc}{$^{\circ}{\rm C}$}

\usepackage{multicol}
\usepackage{cases}

% ---- 合订本 / 上册 / 下册 ----
% “main - nA.tex”“main - nB.tex”先 \def\BookPart{A} 或 {B}，再 \input{main.tex}；
% 直接编译 main.tex 时 \BookPart 未定义，得到合订本。main.tex 中用：
%   \ifPartA ... \fi      合订本和上册包含的内容
%   \ifPartB ... \fi      合订本和下册包含的内容
%   \ifFullBook ... \fi   只有合订本包含的内容
%   \BookPartSuffix       书名后缀：合订本为空，上、下册为“（上）”“（下）”
%   \PartChoice{合订本}{上册}{下册}   按版本取不同的文字（如出版日期）
\newif\ifFullBook
\newif\ifPartA
\newif\ifPartB
\ifdefined\BookPart\else\def\BookPart{}\fi
\def\BookPartNameA{A}
\def\BookPartNameB{B}
\ifx\BookPart\BookPartNameA
  \PartAtrue \def\BookPartSuffix{（上）}
\else\ifx\BookPart\BookPartNameB
  \PartBtrue \def\BookPartSuffix{（下）}
\else
  \FullBooktrue \PartAtrue \PartBtrue \def\BookPartSuffix{}
\fi\fi
\newcommand{\PartChoice}[3]{\ifFullBook#1\else\ifPartA#2\else#3\fi\fi}
